% Bagean: metodhe
% Bab: bukti
% Pérangan: sumber

\documentclass[../../../include/open-logic-section]{subfiles}

\begin{document}

\olfileid{mth}{prf}{res}

\section{Sumber Liyané}

Ana akèh buku babagan cara nyusun bukti ing matématika
kang bisa migunani. Mliginé delengen
\emph{How to Read and do Proofs: An Introduction to
Mathematical Thought Processes} \citep{Solow2013}
lan \emph{How to
Prove It: A Structured Approach} \citep{Velleman2019}.
\href{http://www.people.vcu.edu/~rhammack/BookOfProof/BookOfProof.pdf}{\emph{Book
of Proof}} \citep{Hammack2013} lan
\href{https://scholarworks.gvsu.edu/books/7/}{\emph{Mathematical
Reasoning}} \citep{Sandstrum2019} iku buku babagan
bukti kang bisa diwaca gratis ing internet.
Para filsuf mbokmenawa bakal nemokaké
\emph{More
Precisely: The Math you need to do Philosophy}
\citep{Steinhart2018} minangka pambuka kang
apik tumrap panalaran matématis.

Ana uga manéka pandhuan cekak babagan bukti
ing internet, upamané
\href{https://math.berkeley.edu/~hutching/teach/proofs.pdf}{``Introduction
  to Mathematical Arguments''} \citep{Hutchings2003}
lan
\href{https://eugeniacheng.com/wp-content/uploads/2017/02/cheng-proofguide.pdf}{``How to
  write proofs''} \citep{Cheng2004}.

\subsection{Video Panyemangat}

Rumangsa ora nduwèni semangat kanggo nggarap
tugas? Lagi sedhih? Video-video iki bisa mbantu!

\begin{itemize}
\item \url{https://www.youtube.com/watch?v=ZXsQAXx_ao0}
\item \url{https://www.youtube.com/watch?v=BQ4yd2W50No}
\item \url{https://www.youtube.com/watch?v=StTqXEQ2l-Y}
\end{itemize}

\end{document}
